\subsection{A joint Ferrari test retaining cubic powers and fourth-degree commutators}

\begin{statement}{E.21}{Proved}{Complete compatible corrections in the fourth-degree Ferrari quotient}
Let
\[
 \mathcal A_d=\F_3\langle X,Y\rangle/(\text{monomials of degree}>d),
 \qquad
 K_d=\langle1+X,1+Y\rangle\subset\mathcal A_d^\times.
\]
On the actual free generators $\xi_1,\xi_2$ of the
Ferrari kernel take the automorphisms
\[
 \begin{aligned}
 a:&\quad\xi_1\mapsto\xi_1,
           &&\xi_2\mapsto\xi_2\xi_1,\\
 b:&\quad\xi_1\mapsto\xi_1\xi_2^{-1},
           &&\xi_2\mapsto\xi_2.
 \end{aligned}
 \tag{E.21.1}
\]
Their actual finite image on $K_d$ is denoted $C_d$.
The assignments
\[
 \sigma_1\mapsto(1,a),\qquad
 \sigma_2\mapsto(1,b),\qquad
 \sigma_3\mapsto(\xi_1,a)
 \tag{E.21.2}
\]
define an epimorphism $B_4\to K_d\rtimes C_d$.
Here the cubic powers are retained:
$(1+X)^3=1+X^3$ and $(1+Y)^3=1+Y^3$.

Let $T_4$ be the joint image of (E.21.2) for $d=4$
and the modulus-nine representation $\rho_{4,9}$
of E.19. Its reduction to the E.18 target is
surjective, with elementary abelian kernel
$\F_3^{12}$. Consequently
\[
 |T_4|=344373768,
 \qquad
 |B_4/(\ker(B_4\to T_4)\cap PB_4)|=8264970432.
 \tag{E.21.3}
\]

Form the three-strand core by the six evaluation
maps and all six $B/F_2$ conjugates, exactly as
in (E.19.1), now with target $T_4$, and intersect
it with $K_6$. Denote this kernel by $K_{3,4}$
and put $D_4=B/K_{3,4}$. Its kernel over $G_6$
has dimension seven over $\F_3$, the image of
$K_6\cap F_2'$ has dimension six, and the image
of $N'$ has dimension three.

For each of the 24 retained E.17 $G_6$ points,
the complete gentle norm equations on this
six-dimensional correction image have rank
four; the simultaneous pentagon in the two
targets has rank five. The combined system
is consistent of rank five, hence has an
affine line of corrections. An actual
commutator word in each line is retained.
It satisfies both complete hexagons in
$D_4$, gives a surjective three-strand
shadow, satisfies the full pentagon in
$T_4$, and defines a surjective four-strand
braid map to $T_4$. The three- and
four-strand full twists are preserved.

Thus all 24 CM fibres remain nonempty in
this joint finite test. In the specified
three-strand quotient, $\gamma_3N$ and
$N^{18}$ are killed. It therefore belongs
to the relative class-two coverage of
E.15; the new check is the displayed
four-strand constraint.
\end{statement}
\begin{proof}
The substitutions (E.21.1) are the exact
Ferrari actions of
\cite[(2.9) and (4.15)]{DE_NakamuraTsunogaiYasuda2010}.
Substituting units with constant term one
defines algebra automorphisms of
$\mathcal A_d$: their inverses are the
literal inverse free-group substitutions.
They preserve $K_d$. Direct multiplication
verifies both adjacent braid relations and
the distant braid relation in (E.21.2).
The two literal Ferrari words become
$(\xi_1,1),(\xi_2,1)$, and $a,b$ generate
$C_d$ by its definition. This proves the
epimorphism. It also specifies the source
comparison without replacing the action
of $B_3$ by a homology action.

The calculation first constructs degree
three completely. The actual unit group
$K_3$ has order $2187$, as enumerated
from its two generators. The image
$C_3$ has order $1944$, as enumerated
from the two exact substitution matrices;
its kernel over $\mathrm{SL}_2(\F_3)$
has order $81$. Thus this initial
four-strand target has order $4251528$.
Its joint core with the modulus-nine
target still has correction dimensions
four and three, and all 24 points pass.
These checks are retained in
\texttt{de\_ferrari\_degree3\_model.py}
and \texttt{de\_ferrari\_degree3\_lift.py}.

For degree four the augmentation ideal
has its full 30 monomial coordinates.
Multiplication is ordinary truncated
noncommutative multiplication over
$\F_3$. An automorphism is stored by
the two complete augmentation series
of its generator images. Its linear
operator on all 30 monomials is obtained
by substitution. Composition is exact
matrix substitution; inverse generator
images are obtained degree by degree
from the invertible linear part.
No division by three occurs. The
program \texttt{de\_ferrari\_truncated\_model.py}
checks all Ferrari and braid identities
and compares 1000 successive actual
word reductions, including inverses,
with the independent complete
degree-three model.

Reduction to degree two, followed by
the exact change $z=z_0-u-2v$ of E.18,
gives the specified map to the E.18
target. Its kernel is additive. Indeed
a kernel element of $K_4$ starts in
degree three. An automorphism in the
kernel changes each generator only
in degrees three and four: its
degree-two change is zero by the
E.18 comparison. Applying that change
to a term of degree three first affects
degree five, which is zero here.
Two such automorphism changes also
commute, since their commutator first
affects degree five. Products of two
kernel unit changes first occur in
degree six. Therefore the exact
additive coordinates are the unit's
degree-three and degree-four terms
and the two automorphisms' changes
in those degrees. There are 72
ambient coordinates, with the
actual image computed inside them.
Together with the seven redundant
but injective congruence coordinates
of the modulus-nine factor, these
give 79 ambient linear coordinates
for the joint kernel.

The complete image is obtained by
Schreier generators over the order-648
E.18 quotient. Choose one actual
word representative $t_h$ for each
of its elements. For each of its
three generator edges retain
\[
              t_h\sigma_i t_{h\bar\sigma_i}^{-1}.
 \tag{E.21.4}
\]
These generate the kernel of the
free three-generator group mapping
onto that quotient; their images
therefore generate the full kernel
of the actual joint braid image.
The coordinate matrix has rank
12. Since that kernel is elementary
abelian, this proves the first
order in (E.21.3).

Every braid generator in this
target has order dividing nine,
as directly checked. Its
abelianization, being a quotient
of the cyclic abelianization of
$B_4$, thus has odd order. Every
nontrivial quotient of $\mathfrak S_4$
has a quotient $C_2$; this follows
from its normal subgroups
$1,V_4,A_4,\mathfrak S_4$.
There is consequently no nontrivial
common quotient. Goursat's argument
gives the exact joint product
$T_4\times\mathfrak S_4$ and proves
the second order in (E.21.3).

The induced six-coset representation
of the three-strand core is exactly
(E.19.2), applied in both target
factors. Every $K_6$ image lies
in the just proved additive kernel.
The complete normal image of $N'$
is generated by all 4656 pairwise
commutators of the actual $n_i$,
followed by closure under all $n_i$.
Its rank is three. Adjoining all
sixth powers gives rank seven.
Restricting their exponents to
the full integral rank-95
abelianization kernel of E.19
gives rank six for
$K_6\cap F_2'$. Literal source
words for every basis vector are
retained and checked. Stability
under both $s$ and $r$ is checked,
so these are the exact images
used by the core, including all
commutator and power corrections.

The same norm identities
(E.19.3) and the same literal
gauge identity (E.19.4) apply
to this six-dimensional space.
Their matrices have rank four
in every fibre. Products of
pentagon factors are affine
linear on the abelian correction
kernel, with coefficients given
by conjugation by their actual
base values. The program
\texttt{de\_ferrari\_degree4\_lift.py}
retains the full matrices; the
simultaneous pentagon has rank
five, and the braid-error matrix
has rank zero. The combined
augmented system has rank five
in every case. All basis-pair
substitutions and 100 additional
full substitutions per fibre
check this affine evaluation;
its completeness follows from
the additive-kernel identity,
not from sampling.

An actual correction in each
affine line is then substituted
as a literal word. The program
checks zero integral
$F_2$-abelianization, both
complete hexagons in all
36 induced coordinates, and
the complete pentagon in both
target factors. As in E.19,
generation of the three-strand
core is checked by evaluating
all $n_i$, their commutators
and sixth powers, and closing
under their conjugations. The
kernel image has its full
rank seven in all 24 cases;
the projection to $G_6$ was
already surjective. Hence the
three-strand map is onto.

For each corrected four-strand
triple, the program
\texttt{de\_ferrari\_degree4\_generation.py}
repeats the entire Schreier
calculation (E.21.4). Its
projection still has order
648 and its kernel image has
rank 12, with the same span
as the original target. This
proves generation of $T_4$.
It also checks directly
\[
 (\sigma_1'\sigma_2')^3=(\sigma_1\sigma_2)^3,
 \qquad
 (\sigma_1'\sigma_2'\sigma_3')^4
       =(\sigma_1\sigma_2\sigma_3)^4.
\]

Finally,
\texttt{de\_ferrari\_degree4\_lowercentral.py}
computes the action of every
$n_i$ on the complete rank-three
commutator image. Each action
is the identity, so the image
of $\gamma_3N$ is zero. Every
$n_i^9$ is also trivial in the
induced coordinates. The
commutator image has exponent
three, so class-two collection
shows that every ninth power
in $N$ is trivial there.
Together with the exponent-six
homology component, this proves
$\gamma_3N\,N^{18}\subset K_{3,4}$.
It gives the exact claimed
comparison with E.15 and the
stated finite scope.
\end{proof}
